{\bf Thermal diffusivity of marine sediments} \\
Sediment compaction tends to follow an exponential decay relationship, hence we
can estimate the porosity of the sediment as a function of depth by
\begin{equation}
    \phi(z) = \phi_0 e^{-z/h},
\end{equation}
where $\phi_0$ is the unloaded porosity and $h$ is the scale length for
compaction.
\begin{table}[hbt]
    \centering
    \caption{Properties of oceanic sediments}
    \label{tab:oseds}
    \begin{tabular}{lccc}
        \toprule
        \textbf{Property} &
        \textbf{Seawater} &
        \textbf{Carbonate Ooze} &
        \textbf{Pelagic Clay} \\
        \midrule
        $\phi_0$ & {} & 0.656 & 0.812 \\
        $h$ [m] & {} & 1278 & 664 \\
        $\rho$ [\si{kg.m^{-3}}] & 1030 & 2700 & 2500 \\
        $C_P$ [\si{J.kg^{-1}.K^{-1}}] & 4060 & 815 & 860 \\
        $k$ [\si{W.m^{-1}.K^{-1}}] & 0.6 & 2.24 & 2.3 \\
        \bottomrule
    \end{tabular}
\end{table}
To compute the effective properties, we use the mixing relations,
\begin{align}
    C_P &= C_{P,m} (1 - \phi) + C_{P,w} \phi \\
    \rho &= \rho_m (1 - \phi) + \rho_w \phi \\
    k &= k_m^{(1 - \phi)} k_w^\phi .
\end{align}
For carbonate ooze and pelagic clay, make a plots of
\begin{enumerate}
    \item $\phi(z)$
    \item $k(z)$
    \item $\kappa(z)$

    {\bf Answer:} \\
    It can be a bit of a pain to compute each of these graphs by hand and it is
    therefore best to do so using a spreadsheet or matlab or some other similar
    type program.
    \begin{figure}[hbt]
        \centering
        \includegraphics[]{figures/sedprop.eps}
        \caption{Ocean sediment properties: (a) porosity, (b) thermal
        conductivity, (c) thermal diffusivity for (blue) carbonate and (red)
        pelagic sediments. Note: to compute the thermal diffusivity, you will
        need to find the relationship between thermal diffusivity and
        othermaterial properties, i.e., $\kappa = \frac{k}{\rho C_P}$.}
        \label{fig:sedprop}
        
    \end{figure}
    \item How does compaction affect the rate at which heat can be transferred?

    {\bf Answer:} \\
    Since thermal conductivity and thermal diffusivity both increase with
    decreasing porosity, compaction increases the efficacy of heat transfer.
\end{enumerate}

